% 本文件由 LaTeX 排版 Agent 根据有效 translation.json 自动生成。
% author=Tatian; work_id=0202; title=塔提安致希腊人辞

\chapter{塔提安致希腊人辞}

% structure: p0001 source=h1 target=omitted reason=page-title-already-rendered-by-parent

% structure: p0002 source=h2 target=section reason=relative-heading-rank; base=h2

\section{第一章 希腊人无理声称技艺的发明}

希腊人啊，不要对蛮族怀有如此敌意，也不要以恶眼看待他们的意见。因为你们的制度中有哪一样不是来自蛮族的？最杰出的特耳墨索斯人发明了以梦占卜的技艺；卡里亚人发明了以星宿占卜的技艺；弗里吉亚人和最古老的伊索里亚人发明了以鸟飞占卜的技艺；塞浦路斯人发明了视察牺牲的技艺。你们的天文学归功于巴比伦人；巫术归功于波斯人；几何学归功于埃及人；字母文字的教学归功于腓尼基人。那么，不要再把这些模仿误称为你们自己的发明了。再者，俄耳甫斯教导你们诗歌和歌唱，你们也从他那学习了秘仪。伊特鲁里亚人教导你们雕塑艺术；从埃及人的编年史中你们学会了撰写历史；你们从玛耳叙阿斯和俄林波斯那里获得了吹奏长笛的技艺——这两个弗里吉亚的乡下人编配了牧羊笛的和声。第勒尼安人发明了喇叭；独目巨人们发明了锻造技艺；据赫勒尼科斯告诉我们，一位曾是波斯王后的妇人发明了连接书信木板的方法：她的名字是阿托莎。因此，放下这种自负，不要总是夸耀你们辞藻的优美；因为当你们自我称许时，你们自己的人民自然会附和你们。但一个有见识的人应当等待他人的见证，人们也应当在语言发音上一致。然而，就实际情况而言，唯有你们在普通交谈中也不能说同样的语言；因为多利亚人的说话方式与阿提卡居民不同，伊奥利亚人说话也不像伊奥尼亚人。既然在不应有分歧的地方存在如此大的分歧，我不知该称谁为希腊人。而且，最奇怪的是，你们尊崇并非本土生长的表达，并通过混杂蛮族的词汇使你们的语言变得混杂。因此，我们弃绝了你们的智慧，尽管我一度精通于此；因为，正如那位喜剧诗人所说——\par

\Quote{这些是拾遗的葡萄和闲谈——\newline{}燕子啁啾之处，艺术的败坏者。}\par

然而那些热切追求它的人大声呐喊，像许多乌鸦一样呱呱叫。你们还发明了修辞术来服务于不义和诽谤，出卖你们言论的自由权力以换取报酬，常常在此时将同一件事说成正确，在彼时又说成不好。再者，你们运用诗艺来描述战争、众神的恋情以及灵魂的败坏。\par

% structure: p0006 source=h2 target=section reason=relative-heading-rank; base=h2

\section{第二章 哲学家的恶习与错误}

你们通过追求哲学产生了什么高贵的事物？你们最杰出的人物中，有谁免于虚妄的自夸？第欧根尼，他以他的木桶炫耀独立，却因吃生章鱼而患肠道疾病，因此因贪食丧命。阿里斯提珀，身穿紫袍四处行走，过着放荡的生活，符合他公开宣扬的主张。柏拉图，一位哲学家，因贪食癖好被狄奥尼修斯贩卖。亚里士多德，荒谬地限制天主眷顾，并将幸福置于令人愉悦的事物中，完全违背了他作为导师的职责，谄媚亚历山大，忘记了他不过是个少年；而亚历山大，为了显示他多么精通老师的教诲，当他的朋友不愿崇拜他时，就把他关起来，像熊或豹一样把他带走。他实际上严格遵守了老师的训示，在宴饮中展现勇气与刚毅，用矛刺穿他最亲密、最挚爱的朋友，然后假装悲伤，哭泣并绝食，以免招致朋友的憎恨。我也能嘲笑那些当今遵循他学说的人——他们说月球之下的事物不受天主眷顾；因此，他们比月球更接近大地，位于其轨道之下，自己照看那些未被眷顾的事物；对于那些没有美丽、财富、体力或高贵出身的人，按照亚里士多德的说法，他们没有幸福。让这样的人去哲学思考罢，随我！\par

% structure: p0008 source=h2 target=section reason=relative-heading-rank; base=h2

\section{第三章 对哲学家的嘲弄}

我不能赞同赫拉克利特，他自恃博学而傲慢地说：\Quote{我探索了我自己。}我也不能称赞他将自己的诗篇藏在阿尔忒弥斯神庙中，以便日后作为奥秘发表；那些对此感兴趣的人说，悲剧诗人欧里庇得斯来到那里，阅读了它，逐渐熟记于心，仔细地将赫拉克利特的这一晦涩之作传之后世。然而，死亡显明了此人的愚蠢；因为他患了水肿，虽然他既研究医学又研究哲学，却用牛粪涂满全身，牛粪变硬后，收缩了他整个身体的皮肉，以至于他被撕裂，就这样死去。再者，不能听信芝诺，他声称在终末大火中，同一个人将再次复活，重复以前的行为；例如，安尼图斯和美勒托将控诉，布西里斯将杀害他的客人，赫拉克勒斯将重复他的劳作；在这种终末大火的教义中，他引入了更多恶人而非义人——一个苏格拉底和一个赫拉克勒斯，以及少数同类，但不多，因为恶人将远比善人更多。而且根据他的说法，神性显然将成为邪恶的创造者，居住在阴沟和蠕虫中，以及不敬虔的作恶者中。此外，西西里的火山爆发驳斥了恩培多克勒的空虚自负，因为他虽非神，却几乎假称自己为神。我也嘲笑斐瑞居德的粗陋言论，以及毕达哥拉斯从他继承的教义，还有柏拉图的模仿，尽管有些人持不同看法。谁会赞成克拉特斯的犬儒式婚姻，而不拒绝那些与他相似之人的粗野浮夸言辞，转而探究真正值得关注之事？因此，不要被那些并非哲学家的哲学家的庄严集会所迷惑，他们彼此教条对立，每个人都只是发泄一时的粗俗幻想。此外，他们之间有许多冲突；彼此憎恨；沉溺于相互矛盾的意见，他们的傲慢使他们渴望高位。而且，他们更应不请自来地巴结君王，也不应奉承掌权者，而应等待大人物来找他们。\par

% structure: p0010 source=h2 target=section reason=relative-heading-rank; base=h2

\section{第四章　基督徒只敬拜天主}

希腊人啊，你们为什么想要像拳击手一样，将世俗权力与我们冲突？如果我不愿顺从其中某些人的习俗，为何我要被憎恨如卑鄙的恶徒？君王命令纳税，我乐意缴纳。我的主人命令我如奴隶般服侍，我承认这奴役。人应被当作同伴尊重；唯有天主应被敬畏——祂是肉眼不可见，也不在人类技艺范围之内。只有当祂被否认时，我才不会服从，而是宁愿死也不愿显得虚假和忘恩负义。我们的天主并非在时间中开始存在：唯有祂无始，祂自身就是万有的始。天主是神体，\BibleRef{若 4:24}并非弥漫于物质，而是物质之神的创造者，以及物质中形式的创造者；祂不可见，不可触摸，祂自身就是可感与不可见之物的父。我们从祂的创造认识祂，并通过祂的作为领悟祂不可见的大能。\BibleRef{罗 1:20}我拒绝崇拜祂为我们所造的作品。太阳和月亮是为我们造的：那么，我怎能崇拜我自己的仆役？我怎能称木头石头为神？因为弥漫于物质的神体低于更神圣的神体；即使这与灵魂相似，也不应与完美的天主同等尊敬。也不应将不可言说的天主献上礼物；因为一无所缺的祂不应被我们误解为匮乏。但我要更清楚地阐述我们的观点。\par

% structure: p0012 source=h2 target=section reason=relative-heading-rank; base=h2

\section{第五章　基督徒关于创造世界的教义}

天主在太初；但我们所受的教导是，太初就是圣言的能力。因为宇宙的主宰，祂自身就是一切存有的必然根基（\OriginalTerm{本体}{\GreekText{ὑπόστασις}}），既然尚无任何受造物存在，祂是独一的；但既然祂是全能，自身就是可见与不可见之事的必然根基，万物都与祂同在；与祂同在的，借着圣言的能力（\OriginalTerm{借圣言的能力}{\GreekText{διὰ} \GreekText{λογικῆς} \GreekText{δυνάμεως}}），圣言自身也在祂内，且存于祂内。而借着祂的单纯旨意，圣言发出；圣言并非徒然发出，成为父的首生之工。我们知道祂（圣言）是世界的起始。但祂是因分有而来，而非因切割；因为被切割之物与原本的实体分离，但因分有而来之物，选择了功能，并不使其所取自者有所缺失。就如从一把火炬点燃许多火焰，但第一把火炬的光并不因点燃许多火炬而减少，同样，圣言从父的圣言能力发出，并未使生祂者失去圣言能力。例如，我自己说话，你听见；然而，我说话的人当然不会因传递言语而变得缺乏言语（\OriginalTerm{言}{\GreekText{λόγος}}），而是借着我的声音发声，我努力将你们心中无序的材料整理有序。正如太初所生的圣言，首先为自己创造了必要的物质，进而创造了我们的世界，同样，我仿效圣言，再次被生，并获得了真理，正努力将与我自己同类的混乱材料整理有序。因为物质并不像天主那样无始，也不因无始而与天主同等；它是受生的，不是由任何其他存有产生，而是唯独由万有的创造者使之存在。\par

% structure: p0014 source=h2 target=section reason=relative-heading-rank; base=h2

\section{第六章　基督徒对复活的信德}

因此，我们相信在万物终结之后将有身体的复活；并非如斯多葛派所主张的，依据某些周期的循环，同样的事物毫无目的地产生和毁灭，而是一次性的复活，待我们存在的时期完成之后，并且完全是基于唯有人类所赖以生存的事物结构，以便对他们进行审判。判罚我们的也不是\Person{弥诺斯}或\Person{刺达曼堤斯}，按照神话传说，在他们去世之前没有一个人灵魂受审判；而是造物主，即天主自己，成为仲裁者。尽管你们视我们为无聊的空谈家，但我们并不为此困扰，因为我们对此教义怀有信德。因为正如在我出生之前并不存在，我不知道我是谁，仅仅存在于肉身质料的\OriginalTerm{潜在性}{\GreekText{ὐπόστασις}}中，但出生之后，从先前的虚无状态中，我通过出生获得了对自己存在的确知；同样，既已出生，又因死亡而不再存在，不再被看见，我将再次存在，正如先前我不存在，后来却出生了一样。即使烈火毁灭了我肉身的一切痕迹，世界会接收那蒸发之物；即使被分散在江河海洋，或被野兽撕碎，我仍被储存在一位富有的主的仓库中。虽然穷人和不敬天主者不知道所储存的是什么，但至高无上的天主，一旦祂愿意，必将那唯独祂看得见的实体恢复为其原始状态。\par

% structure: p0016 source=h2 target=section reason=relative-heading-rank; base=h2

\section{第七章 论人的堕落}

天上的圣言，是发自圣父的灵，并且是出自圣言之能的圣言，效法生了祂的圣父，使人成为不朽的肖像，好使正如不朽坏与天主同在，同样，人分享天主的一部分，也能拥有不朽的原则。圣言在创造人之前，也是天使的形成者。这两种受造物各自被造为自由的，可以随己意行动，并不拥有善的本性，这善唯独与天主同在，而是通过人的自由选择得以达至完美，为的是使恶人因自己的过错而堕落，受到公正的惩罚；而义人因自己的德行蒙受应得的赞美，因为在运用自由选择时，他未逾越天主的旨意。这就是关于天使和人的事物构成。圣言的能力本身具有预见未来事件的能力，不是作为命运，而是作为自由主体的选择而发生，它时常预言未来之事的结局；它也通过禁令成为邪恶的禁止者，并成为那些保持良善者的赞美者。当人们依附于一个比其他更狡猾者，因他是首生者而尊崇他，并宣称他是神，尽管他违抗天主的法律，那时圣言的能力便将这愚行的创始者及其追随者排除于自身的一切共融之外。因此，那按天主的肖像受造者，由于更有力的灵离开了他，便成为可朽的；而那首生者因自己的过犯和无知成为邪魔；那些效法他的人，即他的幻象，便成为一群邪魔，并因他们的自由选择而被交付于自己的迷惑中。\par

% structure: p0018 source=h2 target=section reason=relative-heading-rank; base=h2

\section{第八章 邪魔在人类中的罪恶}

但人们形成了他们背叛的\OriginalTerm{材料}{\GreekText{ὑπόθεσις}}。因为，如同掷骰者，向他们展示了星辰位置图之后，他们引入了命运，一种公然的非正义。因为审判者与被审判者由命运安排；杀人者与被杀者、富人与穷人都是同一命运的产物；每一个诞生都被\Person{荷马}所称的那些存在视为一场戏剧表演——\par

\Quote{在诸神中\newline{}升起了不可抑制的笑声。}\par

但那些观看决斗并偏袒一方的人，以及那些结婚、鸡奸、通奸、欢笑、发怒、逃跑、受伤的人，岂不也应被视为凡人吗？因为无论他们通过什么行为向人们显明他们的品格，他们就借此促使听众效法他们的榜样。难道那些邪魔本身，包括以\Person{宙斯}为首的，不也是受命运支配，被与人相同的激情所压倒吗？此外，在那些存在如此大意见分歧者当中，这些存在如何被崇拜呢？因为\Person{瑞亚}，\Place{弗里吉亚山脉}的居民称之为\Person{库柏勒}，因她所爱的\Person{阿提斯}而施行阉割；而\Person{阿佛洛狄忒}喜爱婚姻的结合。\Person{阿耳忒弥斯}是施毒者；\Person{阿波罗}治病。在斩下戈耳工之首后，\Person{波塞冬}的所爱，从中生出飞马\Person{珀伽索斯}和\Person{克律萨俄耳}，\Person{雅典娜}和\Person{阿斯克勒庇俄斯}分了血滴；而他（阿斯克勒庇俄斯）借此救人性命，她（雅典娜）却因同一血成为杀人和战争的煽动者。依我看来，出于对她名誉的考虑，雅典人将她与\Person{赫淮斯托斯}结合所生的儿子归于大地，以免\Person{雅典娜}被认为被\Person{赫淮斯托斯}剥夺了其男性气概，正如\Person{阿塔兰塔}被\Person{墨勒阿革耳}剥夺那样。这个瘸腿的扣环和耳环制造者，很可能，用这些女性装饰品欺骗了那无母的孩子和孤儿。\Person{波塞冬}常去海洋；\Person{阿瑞斯}喜爱战争；\Person{阿波罗}是弹奏基萨拉琴者；\Person{狄俄倪索斯}是\Place{底比斯}人的绝对君主；\Person{克洛诺斯}是弑君者；\Person{宙斯}与自己的女儿交合，她因他怀孕。我还可以引\Place{厄琉息斯}和神秘的龙，以及\Person{俄耳甫斯}，他说——\par

\Quote{关闭大门，不让亵渎者进入！}\par

哀多纽斯拐走了科瑞，他的事迹被制成了秘仪；德默特尔哀悼她的女儿，而某些人则被雅典人欺骗。在勒托之子的神殿区域有一个称为翁法洛斯的地方；但翁法洛斯是狄俄尼索斯的埋葬地。如今我赞美你，哦，达芙妮！——你征服了阿波罗的放纵，证明了他预言能力的虚假；因为，未能预见到将发生在你身上的事，他从他的技艺中一无所获。让那远射之神告诉我，仄费洛斯是如何杀死雅辛托斯的。仄费洛斯战胜了他；正如那位悲剧诗人所说的——\par

\Quote{一阵微风是诸神最尊贵的战车，} —\par

被一阵微风吹拂，阿波罗失去了他的挚爱。\par

% structure: p0026 source=h2 target=section reason=relative-heading-rank; base=h2

\section{第九章 他们引发迷信}

魔鬼就是这样的；正是他们制定了命运的学说。他们的基本原则是将动物置于天上。因为地上的爬行动物、水中的游鱼以及山上的四足兽——他们被逐出天庭后与之同住——这些都被魔鬼赋予天上的尊荣，为的是使他们自己被认为仍留在天上，并通过在那里安置星座，使地上非理性的生活显得合理。因此，意气风发者和被劳作压垮者、节制者和放纵者、贫困者和富有者，他们之所以如此，完全取决于他们出生时的主宰者。因为黄道带的描绘是诸神的工作。而且，当他们中的一个的光芒占优势时，正如他们所说的，它便剥夺了所有其他星辰的尊荣；而此刻被征服的，在另一时间又获得优势。七大行星也乐在其中，仿佛它们是在掷骰子取乐。但我们超越于命运之上，我们不是认识\OriginalTerm{漂泊不定的}{\GreekText{πλανητῶν}}魔鬼，而是认识了唯一一位不漂泊的主；并且，由于我们不追随命运的引导，我们便拒绝它的立法者。我恳求你们告诉我：特里普托勒摩斯在雅典人悲伤之后播种了小麦，成了他们的恩人吗？为什么德默特尔在失去女儿之前不是人类的恩人？厄里戈涅的犬出现在天上，还有阿耳忒弥斯的助手的蝎子、喀戎半人马、被分开的阿尔戈号、卡利斯托的熊。然而，在这些事迹发生之前，诸天是如何没有装饰的呢？而根据某些人的说法，因为\OriginalTerm{宙斯名字的第一个字母}{\GreekText{Διός}}，而将德尔托托姆置于星辰之间，这对谁来说不会显得荒谬呢？为什么萨丁岛和塞浦路斯不在天上受到尊荣？为什么与宙斯分享王国的兄弟们的名字的字母没有被固定在那里呢？克罗诺斯被捆绑并被逐出他的王国，又怎能成为命运的管理者？他自己不再统治，又如何能赐予王国？因此，摒弃这些荒谬之谈，不要因为不公正地仇恨我们而成为违犯者。\par

% structure: p0028 source=h2 target=section reason=relative-heading-rank; base=h2

\section{第十章 嘲笑异教神祇}

有人类变形的传说：在你们那里，神灵也变形了。瑞亚变成一棵树；宙斯为了珀耳塞福涅变成一条龙；法厄同的姐妹们变成了白杨树；勒托变成了一只微不足道的鸟，为了她，如今的提洛岛曾被称为奥尔提基亚。一个神灵，竟然变成了一只天鹅，或一只鹰的形状，并让伽倪墨得斯做他的司酒者，以卑劣的爱情为荣。我怎能敬畏那些贪图礼物、若得不到就发怒的神灵？让他们拥有他们的命运吧！我不愿崇拜漂泊的星辰。贝勒尼基的那一缕头发是什么？在她去世之前，她的星辰在哪里？死去的安提诺乌斯是如何被固定在月亮上，成为一个英俊青年的？是谁把他带到那里？除非，也许就像人们为了报酬而作伪证，当他们嘲弄神灵说国王们升上了天堂时，同样地，有人也把这个人置于众神之中，并以荣耀和报酬作为酬报。你们为什么抢夺天主？你们为什么羞辱祂的创造？你们宰杀一只羊，却又崇拜同一只动物。公牛在天上，你们却宰杀它的形象。跪者座压碎一只毒兽；吞食造人者普罗米修斯的鹰受到崇拜。天鹅确实是高贵的，因为它是一个奸淫者；狄奥斯库里兄弟，轮流活着，抢走琉喀波斯女儿们的人，也是高贵的！海伦更出色，她抛弃了黄发的墨涅拉俄斯，跟随了缠头巾、饰金帕里斯。索弗戎也是一个正义的人，他把这个奸妇运到了福地！但廷达瑞俄斯的女儿并没有被赋予不朽，欧里庇得斯明智地描绘了这个女人被俄瑞斯忒斯杀死的情景。\par

% structure: p0030 source=h2 target=section reason=relative-heading-rank; base=h2

\section{第十一章 人的罪不是由于命运，而是由于自由意志}

既然如此，当我看到这些命运的管理者时，我怎能承认这种按照命运的出生？我不愿成为国王；我不渴求富裕；我拒绝军事指挥；我憎恶淫乱；我不因贪得无厌的获利之心而航海；我不为花冠而竞赛；我摆脱了对名声的狂热渴望；我轻视死亡；我超越一切疾病；忧伤不吞噬我的灵魂。我是奴隶吗？我忍受奴役。我是自由人吗？我不夸耀我的良好出身。我看到同一太阳为所有人照耀，同一死亡为所有人预备，无论他们生活在享乐或匮乏中。富人播种，穷人同享那播种的果实。最富有的人死去，乞丐也有同样的生命极限。富人缺乏许多东西，只有靠众人对他们的评价才显得荣耀；但穷人和欲望适中的人，只寻求适合自己命运的事物，更容易达到目的。你怎么会因贪婪而注定失眠？你怎么会注定常常攫取，又常常死去？向世界死去，拒绝其中的疯狂。向天主而活，借认识祂而脱去旧我。我们不是被造来死的，而是因自己的过错而死。我们的自由意志毁灭了我们；我们本是自由的，却成了奴隶；我们因罪被出卖。天主没有创造任何邪恶；我们亲自显出了邪恶；但既显出了邪恶，我们也能再次拒绝它。\par

% structure: p0032 source=h2 target=section reason=relative-heading-rank; base=h2

\section{第十二章 两种灵体}

我们认识到两种灵体，一种称为\OriginalTerm{灵魂}{\GreekText{ψυχή}}，另一种大于灵魂，是天主的肖像和模样：两者都存在于最初的人中，使之在某方面是\OriginalTerm{物质的}{\GreekText{ὑλικοί}}，在另一方面超越物质。情形如下：我们可以看到，整个世界的结构和一切受造物都由物质产生，而物质本身由天主造成；因此，一方面在物质被分离为部分之前，它可被视为粗糙而未成形的，另一方面在分离之后，它被安排得美丽而有秩序。所以在那个分离中，诸天由物质造成，其中的星辰也是如此；大地及其上的一切都有类似的结构：因此万物有一个共同的起源。但尽管情形如此，由物质造成的物中仍有某些差异，以至于一个更美，另一个虽美却被更好的超越。因为正如身体的结构受到同一管理，并从事那造成其存在的原因的工作，然而尽管如此，身体中有不同的尊荣：眼睛是一回事，耳朵是另一回事，头发的分布和肠道的分配、骨髓与骨骼及肌腱的组合又是另一回事；虽然一部分不同于另一部分，但在它们的安排中却有如同音乐合奏般的和谐——同样，世界按照其造物主的能力，包含一些更灿烂的事物，另一些则不同，并因造物主的意愿而接受了物质的灵体。那些不骄傲地拒绝那些最神圣的、随时间而被记录下来的解释的人，能够分别感知这些事物；这些解释使研习者成为天主的极大热爱者。因此，你们称为魔鬼的那些存在，从物质获得它们的结构，并得到内在于物质的灵体，便变得放纵和贪婪；少数转向更纯洁的事物，但另一些选择了物质中较低劣的部分，并使自己的生活与之一致。这些由物质产生、却远离正确行事的存在，你们希腊人竟敬拜！因为它们因自己的愚妄转向虚荣，并挣脱了权威的缰绳，进而成为天主的劫掠者；万有的主容忍它们嬉戏，直到这世界终结而消散，审判者出现，所有那些受魔鬼攻击却追求认识完美天主的人，在审判的日子将因他们的斗争而获得更完全的见证。因此，在星辰中有一种灵体，在天使中，在植物和水中，在人中，在动物中——但虽是同一灵体，它自身却有差异。我们说的这些事并非道听途说，也非基于可能的猜测和诡辩的推理，而是使用了某种更神圣的言语，你们这些愿意的人，请赶紧来学习。你们若没有轻视西徐亚人\Person{阿那卡雪斯}，就请不要鄙弃遵循一种野蛮法典之人的教导。请以如同接受巴比伦人的预言那样，同样接受我们的教义。请听我们说话，即使如同你们听一棵神谕橡树那样。然而刚才提到的事是癫狂魔鬼的诡计，而我们谆谆教导的教义却远远超越世人的理解。\par

% structure: p0034 source=h2 target=section reason=relative-heading-rank; base=h2

\section{第十三章 灵魂不朽的理论}

灵魂本身并非不朽，希腊人啊，而是可朽的。然而，灵魂也有可能不死。如果它不认识真理，它就死了，并随同身体分解，但最终在世界终穷与身体一同复活，在不朽中接受惩罚性的死亡。但另一方面，如果它获得了对天主的认识，它就不会死，尽管暂时分解。就其本身而言，它是黑暗，其中毫无光明。这话的意义就在于：\Quote{黑暗不能领悟光。}\BibleRef{若 1:5}因为灵魂不保存圣神，而是被圣神保存，光明领悟黑暗。\OriginalTerm{圣言}{Logos}确实是天主的光，而无知的灵魂是黑暗。因此，如果灵魂保持孤独，它就向下趋向物质，并随肉躯死去；但如果它与圣神结合，就不再无助，而上升至圣神引领的地方：因为圣神的居所在上，而灵魂的根源在下。起初，圣神是灵魂的恒常伴侣，但圣神离弃了灵魂，因为灵魂不愿跟随。然而，灵魂保留了一丝圣神的能力，但由于分离而不能辨识完美，在寻求天主时，它在迷途中为自己塑造了许多神，追随魔鬼的诡辩。但天主之神并非与所有人同在，而是寓居于生活正义的人，并亲密地与灵魂结合，藉着预言向其他灵魂宣告隐藏之事。而那些服从智慧的灵魂，吸引了同源的圣神；但悖逆者拒绝了受苦的天主之仆役，显明自己是反抗天主者，而非崇拜祂的人。\par

% structure: p0036 source=h2 target=section reason=relative-heading-rank; base=h2

\section{第十四章 魔鬼将比人受更严厉的惩罚}

你们也是如此，希腊人啊——言辞滔滔不绝，但心思异常乖僻；你们承认众多神的统治，而非独一神的管辖，习惯于追随魔鬼，仿佛它们是有大能的。因为，正如无人性的强盗惯于凭胆量制服同类，魔鬼亦极尽邪恶，完全欺骗了你们中间因无知和虚假表象而放任自流的灵魂。这些存在确实不易死亡，因为它们不分享肉体；但活着时，它们操练死亡之道，每当它们教导追随者犯罪，它们自己也死亡。因此，它们目前的主要特征，即不像人一样死亡，在将要受罚时仍会保留：它们将不会分享永生，从而在蒙福的不朽中接受永生取代死亡。正如我们，现在容易死去，之后会愉悦地接受不朽，或痛苦地接受不朽，同样，魔鬼——它们滥用今生作恶，即使活着也不断死亡——将来也会具有同样的不朽，如同它们活着时所具有的那样，但其性质类似于那些甘愿在生前执行魔鬼命令之人的不朽。难道不是由于人寿命短促，罪恶的种类较少，而魔鬼方面因其无限制的存在，犯罪更为丰富吗？\par

% structure: p0038 source=h2 target=section reason=relative-heading-rank; base=h2

\section{第十五章 与圣神结合的必要性}

然而，我们如今应当寻求我们曾拥有但已失去的，就是将灵魂与圣神结合，并努力追求与天主的合一。人的灵魂由许多部分构成，并非单一；它是复合的，以便通过身体显现自己；因为灵魂没有身体绝不能单独显现，肉躯没有灵魂也不能复活。人并非如呱噪的哲学家所说，仅是理性动物，具备理解力和知识；因为按照他们的话，甚至无理性的动物也显得具有理解力和知识。但唯独人是\OriginalTerm{天主的肖像和模样}{image and likeness of God}；我所说的人，并非指行为与动物相似者，而是指远远超越单纯人性——达到天主本身——的人。我们已在关于动物的论文中更详细地讨论了这个问题。但现在要谈论的主要点是，何谓天主的肖像和模样所指。无可比拟者无非是抽象存在；而被比拟者无非是相似之物。完美的天主无肉躯；而人是肉躯。肉躯的纽带是灵魂；包围灵魂的是肉躯。这就是人的构成性质；如果它像一座圣殿，天主乐意藉着圣神——祂的代表——住在其中；但如果它不是这样的住所，人仅在清晰语言上胜过野兽——在其他方面，他的生活方式与野兽相同，因为他不像天主的模样。但没有一个魔鬼拥有肉躯；它们的结构是属灵的，像火或空气。只有那些被天主之神内住且坚固的人，才容易看见魔鬼的身体，而其他人——我指那些仅有灵魂的人——完全看不见；因为低级者没有能力理解高级者。为此，魔鬼的本性不容悔改；因为它们是物质和邪恶的反映。但物质想要对灵魂行使主权；而依据其自由意志，它们制定了死亡的法律给人；但人在丧失不朽之后，藉着在信德中服从死亡而战胜了死亡；并藉着悔改，他们得到了召唤，正如经上所说：\Quote{因为他们被造略微逊于天使。}并且，每一个被征服之人，若能拒绝带来死亡的条件，就有可能再次征服。那条件是什么，渴望不朽之人不难明白。\par

% structure: p0040 source=h2 target=section reason=relative-heading-rank; base=h2

\section{第十六章 魔鬼虚张声势展示能力}

但是统治人类的\OriginalTerm{魔鬼}{demons}并非人的灵魂；因为死后它们怎能有所行动呢？除非一个人，活着时没有理解和能力，死后却被认为拥有更多的行动能力。但这也是不可能的，正如我们在别处已证明的。而且很难想象，那被身体肢体所阻碍的不朽灵魂，在离开身体后会变得更有智慧。因为魔鬼，因自身的邪恶而对人类疯狂，用各种欺骗性的戏剧表演扭曲他们本已倾向低下的心智，使他们无法兴起走向通往天堂的道路。但对我们而言，世界上的事物并不隐藏，如果那使灵魂不朽的能力眷顾我们，神圣者就很容易被我们领悟。拥有灵魂的人也看见魔鬼，当它们有时向人显现，要么为了被认为是什么，要么像恶意的朋友对敌人那样伤害他们，或为那些相似它们的人提供尊敬它们的机会。因为，如果可能的话，它们无疑会连同其他受造物一起拉下天堂本身。但现在它们绝不能做到，因为它们没有能力；但它们通过低级\OriginalTerm{物质}{matter}与相似自己的物质作战。若有人想征服它们，让他拒绝\OriginalTerm{物质}{matter}。穿上天上圣神的铠甲，他将能保护一切被它环绕的事物。的确，我们内部的物质有疾病和紊乱；但当这类事情发生时，魔鬼将原因归于自己，每当疾病抓住一个人时，它们就靠近他。有时它们自己用一阵愚昧的风暴扰乱身体的状况；但被\OriginalTerm{圣言}{word of God}击打，它们就惊恐地离去，病人便痊愈了。\par

% structure: p0042 source=h2 target=section reason=relative-heading-rank; base=h2

\section{第十七章 他们虚假地向他们的信徒保证健康}

关于\Person{德谟克利特}{Democritus}的同情与反感，我们能说什么呢？除了按照常言，\Place{阿布德拉}{Abdera}人说话是阿布德拉式的啰嗦。然而，正如那位给城市命名的人——据说是\Person{赫拉克勒斯}{Hercules}的朋友——被\Person{狄俄墨得斯}{Diomedes}的马吃掉，同样，那位以\Person{俄斯塔内斯}{Ostanes}自夸的玛吉，在末日将作为永火的燃料被交付。而你们，如果你们不停止嘲笑，将得到与变戏法者相同的惩罚。因此，希腊人啊，请听我说话，我仿佛从高处向你们发言，不要在嘲笑中将你们自己的缺乏理性转嫁给真理的宣告者。一种患病的\OriginalTerm{感受}{\GreekText{πάθος}}不会由一种\OriginalTerm{反感受}{\GreekText{ἀντιπάθεια}}所消除，也不会因在他身上挂小皮护符而治愈疯子。有魔鬼的来访；生病的人、自称恋爱的人、仇恨的人、希望报复的人，都接受它们为帮助者。而它们的运作方式是这样的：正如字母的形式及其组成的线条本身不能指示意思，但人们为自己发明了思想的符号，通过其特定的组合知道字母顺序要表达的意思；同样，各种根茎以及筋腱和骨骼的相互关系本身不能产生任何效果，但它们是魔鬼的堕落所运用的元素性物质，魔鬼确定了每种物质用于什么目的。而当它们看到人们同意被这类事物服务时，它们就接受他们并使他们成为奴隶。然而，服务于通奸如何能体面？激发人们彼此仇恨如何能高贵？或者，将精神错乱的缓解归于物质而不是天主，如何合宜？因为通过它们的技艺，它们使人们偏离对天主的虔诚承认，引导他们信任草药和根茎。然而，天主，如果祂准备了这些事物来实现人们所期望的事，祂就会是恶事的制造者；但祂自己产生了所有具有善性的事物，而魔鬼的放荡将自然产物用于恶的目的，它们所显露的恶的表象来自它们，而非来自完美的天主。因为为什么当我活着时我完全不是恶的，而现在我死了且不能做什么，我的不能运动甚至没有感觉的遗骸，能产生可感知的效果呢？那以最悲惨方式死去的人如何能帮助报复任何人？如果这是可能的，他更可能从自己的敌人那里保卫自己；能够帮助他人，他更可能成为自己的复仇者。\par

% structure: p0044 source=h2 target=section reason=relative-heading-rank; base=h2

\section{第十八章 他们欺骗，而不是治愈}

但医学及其中所包含的一切，都属于同一类的发明。如果有人因信赖物质而被物质治愈，那么因投靠天主之能力而获治愈，岂不更加有效？正如有害的制剂是物质的复合物，治疗性的药物也同样属于同一种性质。然而，倘若我们摒弃那较为卑劣的物质，有些人却往往试图通过将一种坏东西与另一种结合起来而治愈，并利用坏的东西来达到好的目的。但是，正如与强盗同席的人，即使自己不是强盗，也因与他的亲密关系而分担惩罚；同样，那本身不坏却与坏人交往，并因某种自以为的好处而与他们打交道的人，将会被天主审判者因同一目的上的合伙而惩罚。为什么那信赖物质体系的人不愿信赖天主呢？你为何不接近更有能力的主，反而像狗吃草、鹿食蝮蛇、猪吃河蟹、狮吞猴子那样，试图自己治疗自己？你为什么把自然界的事物神化？当你治愈邻人时，为什么被称为恩人？请服从圣言的能力！魔鬼并不治愈，而是用他们的技艺使人成为俘虏。最可敬的犹斯定正确地将他们斥为强盗。因为，正如有些人惯于绑架人，然后以赎金将人归还给朋友；同样，那些被视为神的，侵入某些人的身体，通过梦境制造他们临在的感觉，命令他们到公众面前，在众人眼前，当他们享尽了这个世界的事物后，便从病人身上飞走，并摧毁他们引发的疾病，将人恢复到先前的状态。\par

% structure: p0046 source=h2 target=section reason=relative-heading-rank; base=h2

\section{第十九章  堕落乃魔祟之根本}

但你们这些对此没有知觉的人，应当由我们这些知晓的人来教导；尽管你们自称蔑视死亡，并自认为万事绰绰有余。但这是你们的哲学家极为欠缺的修养，以致其中有些人从罗马皇帝那里每年领取六百金币，却未履行任何有用的服务，只是为了不让他们连蓄着长须也得不到报酬！克雷森，他在大城中筑巢，在逆性之爱（\OriginalTerm{逆性之爱（男色关系）}{\GreekText{παιδεραστία}}）上超越众人，并且极度贪财。然而，这个自称蔑视死亡的人，却如此惧怕死亡，竟试图将死亡作为惩罚加于犹斯定，甚至加于我，因为他借着宣讲真理揭露了那些哲学家是贪食者和骗子。但除了你们之外，他惯常攻击哪位哲学家呢？如果你们赞同我们的教义，说死亡并不可怕，就不要像阿那克萨戈拉那样，出于对声名的疯狂热爱而追求死亡，而要因认识天主而成为蔑视死亡的人。世界的构造是美妙的，但人活在其中的生命是恶劣的；我们可以看到，那些不认识天主的人，如同在庄严集会中受人喝彩。因为占卜是什么？你们为什么受其欺骗？它是你们世俗欲望的仆人。你们想要打仗，便以阿波罗为杀戮的谋士。你们想要强行带走少女，便选一个神作为同谋。你们因自己的过错而生病；正如阿伽门农希望有十位谋士，你们也希望有神与你们同在。某个女人因饮水而癫狂，因乳香的烟雾而失去理智，你们却说她有预言的天赋。阿波罗是预言家和占卜师；在达佛涅的事上，他自欺了。一棵橡树，据说能发出神谕，鸟儿也能发出预兆！如此，你们倒不如动物和植物了！对你们来说，成为占卜用的树枝，或长出鸟的翅膀，倒真是一件美事！那使你们贪财的，也预言你们发财；那鼓动叛乱和战争的，也预言战争得胜。如果你们超越情欲，就会藐视一切世俗之物。不要憎恶我们这些达到了这种境界的人，而要弃绝魔鬼，追随唯一的天主。\BibleQuote{万物是藉着祂造的；凡受造的，没有一样不是藉着祂造的。\BibleRef{若 1:3}}如果自然产物中有毒物，那是因我们的罪过而附加的。我能证明这些事的完全真理；只要你们听从，相信的人将会明白。\par

% structure: p0048 source=h2 target=section reason=relative-heading-rank; base=h2

\section{第二十章  感谢永远归于天主}

即使你藉药物得医治（我姑且以礼相待承认这一点），你仍应将治愈的见证归于天主。因为这世界仍在拖累我们，而我因软弱倾向于物质。因为灵魂的翅膀曾是完美的圣神，但人因罪将之丢弃后，它便如雏鸟般扑腾坠地。离开了天上的共融，它便渴求与下界之物的交往。魔鬼被驱逐到另一居所；最初受造的人类从他们的处所被逐出：前者从天堂被摔下；后者从地上被赶出，然而并非出于这大地，而是出于一种比现今存在者更为优越的秩序。如今，我们渴慕那原初状态，就必须抛开一切阻碍。诸天并非无限，人啊，乃是有限有界的；在其之外是更优越的世界，没有季节更替从之产生各种疾病，而是享有每一适宜的气候，有永恒的白昼和世人无法接近的光明。那些精心描摹大地的人已经尽可能地向人陈述了其不同区域；但由于无法言说超越之物，因为不可能亲身观察，他们便将潮汐的存在归为原因；并说一个海充满杂草，另一个海充满泥土；有些地方被炎热烧灼，另一些地方寒冷冰冻。然而，我们藉先知的教导得知了我们原先未知的事，他们深信天上的圣神将与灵魂一同穿上可朽的衣裳，预言了其他心智所未曾知晓的事。但每一个赤身露体的人都可能获得这衣裳，并回归其古老的亲属关系。\par

% structure: p0050 source=h2 target=section reason=relative-heading-rank; base=h2

\section{第二十一章 基督徒与希腊人关于天主的教义比较}

我们并非做愚昧之事，希腊人啊，也不是讲无稽之谈，当我们宣告天主以人的形状降生时。我呼吁你们这些责难我们的人，将你们的神话记载与我们的叙述相比较。他们说，\Person{雅典娜}为\Person{赫克托耳}的缘故变成了\Person{得伊福玻斯}的形象，不剪发的\Person{福玻斯}为\Person{阿德墨托斯}牧放跛足的牛群，而\Person{宙斯}的配偶化作老妇来到\Person{塞墨勒}面前。但你们既然如此认真地对待这些事情，又如何能嘲笑我们呢？你们的\Person{阿斯克勒庇俄斯}死了，而那位在\Place{特斯匹亚}一夜之间强掠五十个贞女的人，因自我献身于吞噬的火焰而丧命。\Person{普罗米修斯}被锁在\Place{高加索}山上，因他对人类所做的善行而受罚。在你们的说法中，\Person{宙斯}是嫉妒的，他向人隐藏梦境，欲毁灭他们。因此，看看你们自己的纪念物，请赐予我们赞许，即便只是将其视作与你们相似的传说。然而，我们并不行愚妄之事，你们的传说不过是无稽之谈。如果你们谈论诸神的起源，你们也宣称他们是有死的。为什么\Person{赫拉}现在从未怀孕？她老了吗？或者没有人可以告诉你们？现在请相信我，希腊人啊，不要将你们的神话和诸神解释为寓意。如果你们试图这样做，你们所持有的神性就会被你们自己所推翻；因为，如果你们那里的魔鬼就如所说的那样，他们的品性是无价值的；或者，如果被视为自然力量的象征，他们就不是所说的那样。但我不能说服自己去崇拜自然元素，也不能说服我的邻人。\Place{兰普萨库斯}的\Person{墨特罗多洛}在他的\WorkTitle{论荷马著作}中非常愚蠢地论证，将一切都变成寓意。因为他说，\Person{赫拉}、\Person{雅典娜}和\Person{宙斯}都不是那些为他们奉献神圣围地和树林的人所认为的那样，而是自然的部分和元素的某些安排。\Person{赫克托耳}、\Person{阿喀琉斯}、\Person{阿伽门农}以及所有希腊人，和蛮族连同\Person{海伦}和\Person{帕里斯}，既然具有相同的本性，你当然会说他们只是为了诗歌的机制而被引入，这些人物没有一个真实存在。但这些我们只是出于论证的目的而提出；因为甚至不允许将我们的天主观念与那些在物质和泥泞中翻滚的人相比较。\par

% structure: p0052 source=h2 target=section reason=relative-heading-rank; base=h2

\section{第二十二章 对希腊人节庆的嘲笑}

你们的教导又是怎样的呢？谁不会鄙视你们的庄严节庆呢？那些节庆为敬奉邪恶的魔鬼而举行，使人蒙受耻辱。我常常看见一个人——我曾惊愕地看见，而这惊愕又以鄙夷告终，想到他内在是一回事，外表却伪装成不是的东西——过分装出一副优雅的样子，沉溺于各种柔弱；有时眼珠乱转；有时双手挥舞，脸涂泥浆，狂乱不已；有时扮作\Person{阿佛洛狄忒}，有时扮作\Person{阿波罗}；独自一个控告所有神祇，迷信的缩影，诋毁英雄事迹，扮演谋杀，记录奸淫，疯狂的仓库，娼优的教师，死刑的煽动者——然而这样的人却受到所有人的赞美。但我摒弃了他的一切谎言、不敬、行为——简言之，这整个人。但你们却被这样一些人牵着鼻子走，同时却辱骂那些不参与你们活动的人。我没有心思张口呆望一群歌手，也不愿同情一个以不自然方式眨眼和做手势的人。你们中间有什么奇妙或非凡的事呢？他们用矫揉造作的声调说出猥亵之言，做出下流的动作；你们的女儿和儿子在戏台上观看他们教授奸淫的功课。你们的讲堂真是可敬的地方，在那里，夜间犯下的每一桩卑鄙行为都被高声宣告，听众以聆听可耻的演讲为享乐！你们的说谎的诗人也是可敬的，他们通过虚构的谎言引诱听众偏离真理！\par

% structure: p0054 source=h2 target=section reason=relative-heading-rank; base=h2

\section{第二十三章 论拳击手与角斗士}

我曾见过一些人因身体锻炼而疲惫不堪，背负着肉身的重担，在他们面前摆着奖品和花冠，裁判们鼓励他们，不是走向德行的事迹，而是走向暴力与纷争的竞赛；那善于击打者便得了冠冕。这些还是较小的恶；至于更大的恶，谁不避而不谈呢？有些人为了放荡而耽于怠惰，便卖身求死；贫困者出卖自己，而富人则买别人来杀他。为此，见证人坐定，拳击手一对一搏斗，毫无理由，也没有人下到场中去援救。这样的表演难道为你们增光吗？你们中的首领招募一队嗜血的凶手，承诺供养他们；这些暴徒被他派出去，你们则聚集观看这景象，成为裁判，一部分裁判那裁决者的恶行，一部分裁判那参与搏斗之人的恶行。谁错过了这杀人的表演便懊恼，因为他未能成为那邪恶、不虔敬、可憎之事的观众。你们宰杀动物以食其肉，还购买人来为灵魂提供一场人肉宴席，以最不虔敬的流血滋养它。强盗为掠夺而杀人，而富人买角斗士却是为了让他们被杀。\par

% structure: p0056 source=h2 target=section reason=relative-heading-rank; base=h2

\section{第二十四章　论其他公共娱乐}

我从那被\Person{Euripides}搬上舞台、疯狂咆哮、扮演\Person{Alcmæon}弑母的人身上能得到什么好处？他甚至不保持自然的行为，而是张着嘴、手持剑走来走去，尖声叫喊，穿着不合男儿身份的袍子被烧死。此外，\Person{Acusilaus}的神话故事和同类的韵文作家\Person{Menander}也当被弃绝。我为何要敬佩那神话中的吹笛手？我为何要像\Person{Aristoxenus}那样忙于\Place{底比斯}的\Person{Antigenides}？我们将这些无价值之物留给你们；你们要么相信我们的道理，要么像我们一样放弃你们的道理。\par

% structure: p0058 source=h2 target=section reason=relative-heading-rank; base=h2

\section{第二十五章　哲学家的夸耀与争吵}

你们的哲学家成就了什么伟大奇妙的事？他们露着一侧肩膀；他们留长发；他们蓄胡须；他们的指甲如同野兽的爪子。虽说他们一无所缺，却像\Person{Proteus}一样，需要皮匠做他们的行囊，织工做他们的外衣，樵夫做他们的手杖，还需要富人和厨师来满足他们的贪食。\newline{} 与狗竞争的人哪，你不认识天主，于是转向模仿无理性的动物。你以专横的姿态当众呼喊，并自行报复；若你一无所获，便肆意谩骂；哲学对你而言成了赚钱的技艺。你追随\Person{Plato}的学说，而一个\Person{Epicurus}的门徒起来反对你。你又想成为\Person{Aristotle}的门徒，而一个\Person{Democritus}的信徒便讥讽你。\Person{Pythagoras}说他是\Person{Euphorbus}，他继承了\Person{Pherecydes}的学说；但\Person{Aristotle}却攻击灵魂的不朽。你们从前辈那里接受相互冲突的学说，你们这些不和谐之人，正与和谐者作战。你们中一人说天主是有形的，我却说祂是无形的；说世界是不可毁灭的，我却说它将要被毁灭；说每隔一段时间会有大火，我却说它只发生一次；说\Person{Minos}和\Person{Rhadamanthus}是审判者，我却说天主自己就是审判者；说唯独灵魂享有不朽，我却说肉身也享有不朽。\newline{} 我们损害了你们什么，\Place{希腊}人啊？你们为何憎恨那些跟随天主圣言的人，仿佛他们是人类中最卑劣的？不是我们吃人肉——你们中编造此事的人是收买的伪证人；在你们中，\Person{Pelops}虽为\Person{Poseidon}所爱，却被做成给诸神的晚餐；\Person{Kronos}吞食自己的儿女；\Person{Zeus}吞下\Person{Metis}。\par

% structure: p0060 source=h2 target=section reason=relative-heading-rank; base=h2

\section{第二十六章　对希腊学术的嘲讽}

不要再炫耀那些从别人那里得来的言论，不要像乌鸦一样借别人的羽毛来装饰自己。如果每个城邦都收回它对你言辞的贡献，你的谬论就会失去力量。在探究天主是什么时，你却对自身无知；仰视天空目瞪口呆，却陷入陷阱。读你们的书犹如走进迷宫，读者就像达那伊得斯的桶。你们为什么要划分时间，说一部分是过去，一部分是现在，另一部分是未来？当现在存在时，未来如何能流逝？就像航行中无知的人，当船向前行驶时，以为山在移动，你们不知道是自己在流逝，而\OriginalTerm{时间}{\GreekText{ὁ} \GreekText{αἰών}} 只要造物主愿意它存在，就始终是现在。为什么我要为表达自己的意见而被追究责任？你们为什么如此急于全部否定它们？你们不是和我们一样出生，处于同一世界的治理之下吗？为什么说智慧只在你们那里？你们并没有另一个太阳，没有另一种星辰的升起，没有更优越的起源，没有比其他人更优越的死亡。语法学家是这些空谈的起源；你们这些分派智慧的人，与合乎真理的智慧隔绝，将各个部分的名称分派给特定的人；你们不认识天主，在激烈的争执中互相毁灭。因此，你们全都一文不值。你们自以为是讨论的唯一权利，却像盲人与聋子交谈。为什么你们不懂建造，却摆弄建造的工具？为什么你们只忙忙碌于言辞，却远离行动，因赞美而自大，因不幸而沮丧？你们的行事方式违背理性，公开炫耀，却将教义藏在角落里。我们发现你们就是这样的人，便离开了你们，不再关心你们的教义，而是跟随天主的圣言。人啊，你为什么让字母相互争斗？为什么像拳击一样，用你矫揉造作的雅典腔调让声音碰撞？你本应更自然地说话。如果你采用雅典方言，却并非雅典人，请问为什么不像多里安人那样说话？为什么一种显得更粗野，另一种更便于交往？\par

% structure: p0062 source=h2 target=section reason=relative-heading-rank; base=h2

\section{第二十七章 基督徒被不公正地仇恨}

如果你们坚持\Emph{他们}的教义，为什么又因我选择我所赞同的教义观点而攻击我？强盗不会因他所拥有的名字而受惩罚，只有事实真相查清后才被定罪，而我们却未经审查就遭到谩骂，这难道不合理吗？狄亚戈拉斯是雅典人，你们却因他泄露了雅典的秘仪而惩罚他；然而你们读了他的弗里吉亚著作，却仇恨我们。你们拥有利奥的注释，却对我们驳斥它们感到不悦；你们手中有阿皮翁关于埃及诸神的意见，却谴责我们为最不敬神的人。奥林匹亚宙斯的坟墓在你们中间被展示，尽管有人说克里特人是说谎者。你们众多神的集会毫无价值。虽然轻视诸神的伊壁鸠鲁担任火炬手，但我绝不会因此向统治者隐瞒我对天主及祂治理宇宙的看法。为什么你们劝我违背自己的原则？你们说你们轻视死亡，为什么却劝我们用技巧逃避死亡？我没有鹿的心；但你们对辩证法的热情却类似于忒耳西忒斯的饶舌。我怎能相信有人告诉我太阳是一块炽热的铁、月亮是一片土地？这种断言只是言辞之争，并非对真理的清醒阐述。相信希罗多德关于赫拉克勒斯史事的记载，说有一只狮子从上层大地下来并被赫拉克勒斯杀死，这不是愚蠢吗？还有阿提卡风格、哲学家的复杂推论、三段论的外表、地球的测量、星辰的位置、太阳的轨道，这有什么用处？致力于这些研究，就像是把意见当作法律强加给自己的人的工作。\par

% structure: p0064 source=h2 target=section reason=relative-heading-rank; base=h2

\section{第二十八章 谴责希腊人的律法}

因此，我也拒绝你们的律法；因为应当有一个全人类共同的政体；但现在有多少城邦就有多少不同的法典，以致在有些地方被视为可耻的事，在另一些地方却是荣誉的。希腊人认为与母亲发生关系是不合法的，但波斯的玛各却将此视为最恰当的；鸡奸被蛮族谴责，但罗马人却赋予它某些特权，因为他们像牧马一样试图收集男孩。\par

% structure: p0066 source=h2 target=section reason=relative-heading-rank; base=h2

\section{第二十九章 塔提安皈依的叙述}

因此，眼见这些事，并已参加奥秘仪式，且处处考察由羸弱者和被鸡奸者行进的宗教仪式，我发现罗马人的拉提亚里斯·朱庇特喜爱人血和被杀之人的血，而阿尔特弥斯在大城附近也批准同样行为，此处的恶鬼和彼处的恶鬼煽动行恶——我退隐独处，寻求如何能发现真理。当我全神贯注于此问题时，我偶然遇到某些野蛮人的著作，其古老无与伦比于希腊人的意见，其神圣无与伦比于他们的错误；我被那不加修饰的语言、作者朴拙的笔法、对未来事件的先见、卓越的训诫，以及宇宙的统治集中于独一之神的宣告所说服，从而信靠这些著作。我的灵魂受天主教导，我看出前一类著作引人定罪，但这些却终结世上的奴役，将我们从众多统治者和无数的暴君手中解救出来，所赐予的并非我们先前未曾领受的，而是我们已领受但因错误而未能保持的。\par

% structure: p0068 source=h2 target=section reason=relative-heading-rank; base=h2

\section{第三十章 他如何决心抵抗魔鬼}

因此，受启蒙并被教导这些事之后，我愿像抛弃幼稚的愚行一样抛弃先前的错误。因为我们知道邪恶的本性就像最小的种子；它从微小的开端变得强大，但如果我们听从天主的言语并不放纵自己，它又将再次被毁灭。因为祂藉着某种\Quote{隐藏的宝藏}成为我们一切所有的主宰；挖掘这宝藏时我们确实满身尘土，但我们却确保它为我们固定的财产。接受这全部宝藏的人已获得了最宝贵财富的支配权。就让这些话对我们的朋友们说。但你们希腊人，我能说什么呢，除非请求你们不要辱骂比你们更好的人，也不要因为他们被称为野蛮人而以此作嘲弄的借口。因为如果你们愿意，你们就能发现人不能理解彼此语言的原因；因为对于那些愿意考察我们教义的人，我将对这些教义作一个简单而详尽的说明。\par

% structure: p0070 source=h2 target=section reason=relative-heading-rank; base=h2

\section{第三十一章 基督徒的哲学比希腊人的更古老}

但现在似乎是时候证明我们的哲学比希腊人的体系更古老了。让\Person{梅瑟}和\Person{荷马}作为我们的界限，二者都极为古老；前者是最古老的诗人兼史学家，后者是一切野蛮人智慧的奠基人。那么，让我们在他们之间作比较；我们将发现我们的教义不仅比希腊人的更古老，而且甚至先于文字的发明。我将不从我们内部引证见证人，而宁愿诉诸希腊人。这样做前者将是愚蠢的，因为你们不会允许；但后者会使你们惊讶，因为当我用你们自己的武器与你们交锋时，我提出的论点你们从未怀疑。现在，最古老的作家们已经探究过荷马的诗歌、他的家世和他生活的时代——\Person{雷吉翁的特阿革涅斯}，他生活在冈比西斯时代；\Person{塔索斯的斯特西姆布罗托斯}和\Person{科洛丰的安提马科斯}；\Person{哈利卡纳苏斯的希罗多德}；\Person{奥林图斯的狄奥尼修斯}；此后，\Person{库迈的埃福罗斯}、\Person{雅典的菲洛克罗斯}、\Person{墨伽克勒斯}和\Person{逍遥派的卡梅莱翁}；后来是文法学家\Person{泽诺多托斯}、\Person{阿里斯托芬}、\Person{卡利马科斯}、\Person{克拉特斯}、\Person{埃拉托斯特尼}、\Person{阿里斯塔克斯}和\Person{阿波罗多罗斯}。其中，\Person{克拉特斯}说他生活在赫拉克勒斯后裔归来之前，特洛伊战争之后八十年内；\Person{埃拉托斯特尼}说是在伊利昂陷落后第一百年之后；\Person{阿里斯塔克斯}说大约在伊奥尼亚移民时期，即那事件之后一百四十年；但根据\Person{菲洛克罗斯}，在伊奥尼亚移民之后，雅典的阿基普斯执政官任期内，特洛伊战争后一百八十年；\Person{阿波罗多罗斯}说是在伊奥尼亚移民后一百年，那将是特洛伊战争后二百四十年。有人说他生活在奥林匹亚赛会之前九十年，那将是特洛伊陷落后三百一十七年。其他人则将其推迟到更晚的时期，说\Person{荷马}与\Person{阿尔基洛科斯}同时代；但\Person{阿尔基洛科斯}生活在第二十三届奥林匹亚赛会期间，\Person{吕底亚的古格斯}时代，特洛伊之后五百年。因此，关于前述诗人（我指\Person{荷马}）的时代以及那些谈到他的人们之间的分歧，我们已以概要方式为那些能够准确探究的人说了足够的话。因为有可能表明，关于事实本身的意见也是错误的。因为，如果指定的日期不一致，历史就不可能是真实的。因为写作中的错误原因是什么，岂非叙述不真实之事吗？\par

% structure: p0072 source=h2 target=section reason=relative-heading-rank; base=h2

\section{第三十二章 基督徒的教义反对纷争，适合所有人}

但在我们中间，并无虚荣的欲望，也不沉溺于纷歧的见解。因为我们弃绝了世俗的流行观念，服从天主的命令，遵行永生圣父的律法，拒绝一切基于人意的学说。我们中间的富人追求我们的哲学，穷人也无偿地领受教诲；因为出自天主的事物超越世俗礼物的回报。因此我们接纳所有愿意聆听的人，甚至老妇和少年；简言之，各种年龄的人都受到我们的尊敬，而各种放荡行为则被远远摒弃。我们说话时不说谎言。倘若你们的不信能受到遏制，那将是一件好事；但无论如何，让我们的主张由天主的审判所确认而坚定。你们若愿意，尽管嘲笑；但日后必当哀哭。难道这不荒谬吗？你们说，\Person{涅斯托耳}因年老力衰，解马缰迟缓，却因试图与年轻人争战而受钦佩；而你们却嘲笑我们中那些与衰老抗争、专注于天主之事的人。你们告诉我们曾有\Person{阿玛宗人}、\Person{塞弥拉弥斯}以及其他若干好战女子存在，却指责我们的少女，谁不会发笑？\Person{阿喀琉斯}是个青年，却被认为胸怀非常博大；\Person{涅俄普托勒摩斯}更年轻，却强壮；\Person{菲罗克忒忒斯}体弱，但神明需要他去攻打\Place{特洛伊}。\Person{忒耳西忒斯}是何许人？他却在军中担任指挥，若不是因愚钝而舌无遮拦，就不会因尖头秃顶而受责难。至于那些想学习我们哲学的人，我们不凭外貌考验他们，也不凭外表判断来者；因为我们主张，尽管身体软弱，人人皆可有思想的力量。但你们的行为却充满嫉妒和十足的愚蠢。\par

% structure: p0074 source=h2 target=section reason=relative-heading-rank; base=h2

\section{第33章 为基督徒女性辩护}

因此，我愿从你们视为尊贵的事上证明，我们的规矩以克制为特色，而你们的规矩则与疯狂相近。你们说我们在女人和男孩、少女和老妇当中胡言乱语，并因我们不与你们同列而嘲笑我们，请听听希腊人中有何等愚昧。因为他们的艺术品献给无价值的对象，而你们对这些艺术品的崇敬甚至超过你们的神祇；你们在有关女人之事上举止不当。\Person{吕西波斯}铸造了\Person{普拉克西拉}的雕像，其诗歌毫无益处；\Person{墨涅斯特拉托斯}铸造了\Person{勒阿耳基斯}的雕像，\Person{塞拉尼翁}铸造了妓女\Person{萨福}的雕像，\Person{瑙基得斯}铸造了\Person{厄里那}（勒斯比亚人）的雕像，\Person{波伊斯库斯}铸造了\Person{密耳提斯}的雕像，\Person{刻菲索多托斯}铸造了拜占庭的\Person{密戎}的雕像，\Person{龚福斯}铸造了\Person{普拉克西戈里斯}的雕像，\Person{安菲斯特拉托斯}铸造了\Person{克利托}的雕像。关于\Person{阿倪塔}、\Person{忒勒西拉}和\Person{密斯提斯}，我该说什么呢？\Person{欧提克拉底}和\Person{刻菲索多托斯}为前者造像，\Person{尼克拉托斯}为第二位，\Person{阿里斯托多托斯}为第三位；\Person{欧提克拉底}为以弗所人\Person{姆涅西亚尔基斯}造像，\Person{塞拉尼翁}为\Person{科里娜}造像，\Person{欧提克拉底}为阿尔戈斯人\Person{塔拉耳基斯}造像。我提及这些女人，为要使你们不把我们中间的事视为异样，并透过眼前所见的雕像，不轻蔑我们中追求哲学的女子。这位\Person{萨福}是一个淫荡、恋慕的女子，歌唱自己的放荡；但我们所有的女子都贞洁，少女们在纺锤旁歌唱神圣的事，比你们的那个姑娘更高贵。因此，你们这些自称是女子门徒的人，却嘲笑我们中持守教义的女子以及她们所参加的庄严集会，当觉羞愧。\Person{格劳喀珀}为你们生下一个何等高贵的孩子？她生下一个畸形儿，正如雅典人\Person{欧克特蒙}之子\Person{尼克拉托斯}所铸的雕像所示。但是，若\Person{格劳喀珀}生下一只象，难道她就该享有公众荣誉吗？\Person{普拉克西特勒斯}和\Person{赫罗多特斯}为你们铸造了妓女\Person{弗里涅}的雕像，\Person{欧提克拉底}铸造了\Person{潘透基斯}的铜像，她因一个嫖客而怀孕；\Person{狄诺墨涅斯}因为派奥尼亚人的王后\Person{贝桑提斯}生下一个黑婴，便费心用他的艺术留存她的记忆。我也谴责\Person{毕达哥拉斯}，他制作了骑在公牛上的\Person{欧罗巴}的雕像；也谴责你们，因艺术技巧而尊崇宙斯的指控者。我嘲笑\Person{弥戎}的技艺，他制作了一头母牛，牛上有一位胜利女神，因为（公牛）劫走\Person{阿革诺耳}的女儿，赢得了奸淫和放荡的奖赏。奥林索斯人\Person{赫罗多特斯}铸造了妓女\Person{格里克拉}和竖琴女\Person{阿尔革亚}的雕像。\Person{布里阿克西斯}铸造了\Person{帕西法厄}的雕像；通过留存她的放荡记忆，几乎可以看出你们的愿望是让当代女子像她一样。一位名叫\Person{墨拉尼珀}的智慧女子，为此\Person{吕西斯特拉托斯}为她造像。但是，你们确实不会相信，我们当中也有智慧女子！\par

% structure: p0076 source=h2 target=section reason=relative-heading-rank; base=h2

\section{第34章 对希腊人所立雕像的嘲弄}

那位吞食婴儿的暴君\Person{法拉里斯}当然配享极崇高的荣誉，而安布拉基亚人\Person{波吕斯特拉托斯}的工艺至今仍将他展示为一位非凡的人物！阿克拉加斯人因他的食人习性而不敢直视他的面容，但如今有教养的人却以能瞻仰其雕像为荣！你们这些观看\Person{波吕尼刻斯}和\Person{厄忒俄克勒斯}雕像的人，竟尊崇亲手足相残者，而不愿与他们的制造者\Person{毕达哥拉斯}一同埋葬这些雕像，岂不羞耻？毁掉这些邪恶的纪念物吧！我何必仅仅为了雕刻家\Person{佩里克吕墨诺斯}而赞叹那位生了三十个孩子之女的雕像？人们应当厌恶地避开那承受极大荒淫果实的人，罗马人将她比作母猪，据说这母猪也因相似原因被视为配享神秘崇拜。\Person{阿瑞斯}与\Person{阿佛洛狄忒}通奸，\Person{安德隆}为他们的孩子\Person{哈耳摩尼亚}制作了雕像。\Person{索弗隆}记录琐碎荒谬之事，却因铸造技艺而更加闻名，其作品至今犹存。不仅他的故事使寓言作家\Person{伊索}永垂不朽，\Person{阿里斯托得摩斯}的雕塑艺术也增加了他的声誉。你们拥有如此多的女诗人，作品纯属糟粕，还有无数的妓女与无价值之男子，却为何不羞于诽谤我们女性的名誉？我何需知道\Person{欧安忒}在逍遥学园生下一子，或因\Person{卡利斯特拉托斯}的艺术而张口惊叹，或凝视\Person{卡利阿得斯}的\Person{涅埃拉}？她是个妓女。\Person{拉伊丝}是娼妓，\Person{图尔努斯}为她立了娼妓纪念碑。你们为何不以\Person{赫费斯提翁}的淫乱为耻，尽管\Person{斐洛}已极其艺术地描绘了他？又因何故尊崇\Person{莱奥卡雷斯}的阴阳人\Person{伽倪墨得斯}，仿佛拥有什么值得赞赏之物？\Person{普拉克西特列斯}甚至为一位不洁的女人制作了雕像。你们理应摒弃这一切，寻求真正值得关注的事物，而不应在接纳\Person{斐莱尼斯}和\Person{厄勒番提斯}的可耻作品的同时，厌弃我们的生活方式。\par

% structure: p0078 source=h2 target=section reason=relative-heading-rank; base=h2

\section{第三十五章 \Person{塔提安}作为目击者发言}

我向你们陈述的这些事并非道听途说。我游历了许多地方，像你们一样研习修辞，接触了许多技艺与发明，最后在\Place{罗马}城寄居期间，我检视了你们带去的众多雕像。因为我不像许多人那样，惯于用他人观点来加强自己见解，而是希望将我亲眼所见、亲身所感清楚告诉你们。因此，我告别了罗马人的傲慢、雅典人的空谈以及他们所有支离破碎的见解，接受了我们蛮族的哲学。我曾开始说明这哲学比你们的制度更古老，但为了讨论一项更迫切的问题而中途搁置；如今是时候尝试谈论其教义了。不要因我们的教导而恼怒，也不要进行一场充满琐碎与秽语的详细反驳，说：\Quote{\Person{塔提安}，自诩超越希腊人，超越无数哲学研究者，开辟了一条新路，接受了蛮族的教义。}因为明明无知的人被一个与他们本性相同的人说服，这有什么可抱怨的呢？或者，根据你们自己的智者所言，活到老学到老，这有什么不合理呢？\par

% structure: p0080 source=h2 target=section reason=relative-heading-rank; base=h2

\section{第三十六章 加色丁人对\Person{梅瑟}古老性的见证}

但就算\Person{荷马}不晚于\Place{特洛伊}战争，就算他与战争同时代，甚至是在\Person{阿伽门农}军中，或者如有人乐意认为的，他生于文字发明之前，那位先前提到的\Person{梅瑟}将被证明比\Place{特洛伊}失陷早了许多年，比\Place{特洛伊}城、\Person{特罗斯}和\Person{达耳达诺斯}的建筑更古老。为证明这一点，我将邀请加色丁人、腓尼基人和埃及人作为证人。我还需多言吗？因为一个声称要使听者信服的人，理应使其叙述极其简洁。巴比伦人\Person{贝罗索斯}，是他们神祇\Person{彼勒}的祭司，生于\Person{亚历山大}时期，为三世之后的\Person{安提约古}撰写了三卷加色丁人史；叙述诸王事迹时，他提到一位名叫\Person{拿步高}的王，曾与腓尼基人和犹太人作战——我们知道这些事由我们的先知预言，并发生在远晚于\Person{梅瑟}的时代，在波斯帝国之前七十年。\Person{贝罗索斯}是非常可信的人，\Person{犹巴}为此作证，他写亚述史时说，他从\Person{贝罗索斯}学到历史；他有两卷关于亚述人的著作。\par

% structure: p0082 source=h2 target=section reason=relative-heading-rank; base=h2

\section{第三十七章 腓尼基人的见证}

继加色丁人之后，腓尼基人的见证如下。他们中有三人：\Person{特奥多图斯}、\Person{许普西克拉特斯}和\Person{摩库斯}；\Person{柴图斯}将他们的著作译成希腊文，并精确撰写了哲学家的传记。前述作者的历史中显示，\Person{欧罗巴}被劫发生在一王在位时，并记述了\Person{默涅拉俄斯}来到腓尼基，以及与\Person{希兰}有关的事，\Person{希兰}将女儿嫁给犹太王\Person{撒罗满}，并为圣殿的建造提供了各种木材。\Place{帕加马}的\Person{默南德尔}也撰写了关于同样事件的历史。但\Person{希兰}的时代大约在\Place{特洛伊}战争时期；而与\Person{希兰}同时代的\Person{撒罗满}，则远晚于\Person{梅瑟}的时代。\par

% structure: p0084 source=h2 target=section reason=relative-heading-rank; base=h2

\section{第三十八章 埃及人将\Person{梅瑟}置于\Person{伊那科斯}统治时期}

\Place{埃及}人也有精确的编年史。\Place{门德斯}的祭司\Person{托勒密}——并非那位国王——是埃及事务的诠释者。这位作者叙述诸王的事迹，说犹太人离开埃及前往所去之地的这件事，发生在\Person{阿摩西斯}王时代，由\Person{梅瑟}领导。他这样说：\Quote{\Person{阿摩西斯}生活在\Person{伊那科斯}王时代。}在他之后，最受尊崇的语法学家\Person{阿皮翁}，在其\WorkTitle{埃及志}第四卷（他共有五卷）中，除了许多其他事情外，还提到\Person{阿摩西斯}在\Place{阿尔戈斯}的\Person{伊那科斯}时代摧毁了\Place{阿瓦里斯}，正如\Place{门德斯}的\Person{托勒密}在他的编年史中所记载的。但从\Person{伊那科斯}到\Place{特洛伊}陷落之间，相隔二十代。论证的步骤如下：——\par

% structure: p0086 source=h2 target=section reason=relative-heading-rank; base=h2

\section{第三十九章 \Place{阿尔戈斯}诸王名录}

\Place{阿尔戈斯}诸王如下：\Person{伊那科斯}、\Person{福洛纽斯}、\Person{阿庇斯}、\Person{克里阿西斯}、\Person{特里俄帕斯}、\Person{阿耳革乌斯}、\Person{福尔巴斯}、\Person{克罗托帕斯}、\Person{斯忒涅拉俄斯}、\Person{达那俄斯}、\Person{林叩斯}、\Person{普洛托斯}、\Person{阿巴斯}、\Person{阿克里西俄斯}、\Person{珀耳修斯}、\Person{斯忒涅拉俄斯}、\Person{欧律斯透斯}、\Person{阿特柔斯}、\Person{提厄斯忒斯}和\Person{阿伽门农}，在他统治的第十八年，\Place{特洛伊}被攻陷。每一个有智慧的人都会非常仔细地观察到，根据\Place{希腊}人的传统，他们没有任何历史著作；因为教授他们文字的\Person{卡德摩斯}是在许多代之后才来到\Place{玻俄提亚}的。但在\Person{伊那科斯}之后，在\Person{福洛纽斯}统治下，他们野蛮游牧的生活才勉强得到遏制，并进入了新的秩序。因此，如果\Person{梅瑟}被证明与\Person{伊那科斯}同时代，他就比\Place{特洛伊}战争早四百年。这一点从\Place{阿提刻}［以及\Place{马其顿}、\Place{托勒密}和\Place{安提阿}］诸王的世系中得到证明。因此，如果\Place{希腊}人中最辉煌的事迹是在\Person{伊那科斯}之后才被记录和知晓，那么显然这些必定是在\Person{梅瑟}之后。在\Person{伊那科斯}之后的\Person{福洛纽斯}时代，\Place{雅典}人中提到\Person{奥古革斯}，在他那时发生了第一次洪水；在\Person{福尔巴斯}时代是\Person{阿克泰俄斯}，\Place{阿提刻}因而被称为\Place{阿克泰亚}；在\Person{特里俄帕斯}时代是\Person{普罗米修斯}、\Person{厄庇墨透斯}、\Person{阿特拉斯}、双性的\Person{刻克洛普斯}和\Person{伊俄}；在\Person{克罗托帕斯}时代发生了\Person{法厄同}的焚烧和\Person{丢卡利翁}的洪水；在\Person{斯忒涅拉俄斯}时代是\Person{安菲克堤翁}的统治、\Person{达那俄斯}来到\Place{伯罗奔尼撒}、\Person{达尔达诺斯}建立\Place{达尔达尼亚}，以及\Person{欧罗巴}从\Place{腓尼基}返回\Place{克里特}；在\Person{林叩斯}时代是\Person{科瑞}的劫掠、\Place{厄琉息斯}神庙的建立、\Person{特里普托勒摩斯}的农耕，以及\Person{卡德摩斯}来到\Place{忒拜}和\Person{米诺斯}的统治；在\Person{普洛托斯}时代是\Person{欧摩尔波斯}对\Place{雅典}人的战争；在\Person{阿克里西俄斯}时代是\Person{珀罗普斯}从\Place{弗里吉亚}迁来、\Person{伊翁}来到\Place{雅典}、第二\Person{刻克洛普斯}、\Person{珀耳修斯}和\Person{狄俄倪索斯}的功业，以及\Person{俄耳甫斯}的弟子\Person{穆赛俄斯}；在\Person{阿伽门农}统治时期，\Place{特洛伊}被攻陷。\par

% structure: p0088 source=h2 target=section reason=relative-heading-rank; base=h2

\section{第四十章 \Person{梅瑟}比异教英雄更古老、更可信}

因此，从前面所说可知，\Person{梅瑟}比古代英雄、战争和魔鬼更古老。我们更应相信年代更早的他，而非\Place{希腊}人，他们无意中从他的学说中汲取，如同从泉源。因为他们中的许多智者，受好奇心驱使，试图篡改他们从\Person{梅瑟}以及与他一样思考哲学的人那里学到的一切：首先，为了被视为有自己独特的东西；其次，为了用某种修辞技巧掩盖那些他们不理解的事情，从而将真理曲解为寓言。但\Place{希腊}学者关于我们的政体和律法历史所说的，以及有多少人和什么样的人写过这些，将在针对那些谈论神圣事物之人的论著中展示。\par

% structure: p0090 source=h2 target=section reason=relative-heading-rank; base=h2

\section{第四十一章}

但最重要的事情是尽力精确地阐明，\Person{梅瑟}不仅比\Person{荷马}更古老，而且比所有在他之前的作家都更古老——比\Person{利诺斯}、\Person{菲拉蒙}、\Person{塔米里斯}、\Person{安菲翁}、\Person{穆赛俄斯}、\Person{俄耳甫斯}、\Person{德摩多科斯}、\Person{斐弥俄斯}、\Person{西彼拉}、克里特的\Person{厄庇墨尼得斯}（他曾来到\Place{斯巴达}）、\Place{普洛孔涅索斯}的\Person{阿里斯泰俄斯}（他写了\WorkTitle{阿里马斯皮亚}）、\Person{马人阿斯玻罗斯}、\Person{伊萨提斯}、\Person{德律蒙}、\Place{塞浦路斯}的\Person{欧克洛斯}、\Place{萨摩斯}的\Person{霍鲁斯}和\Place{雅典}的\Person{普洛那庇斯}更古老。现在，\Person{利诺斯}是\Person{赫拉克勒斯}的老师，而\Person{赫拉克勒斯}比\Place{特洛伊}战争早一代；这从他儿子\Person{特勒波勒摩斯}（他曾在\Place{特洛伊}军中服役）可以清楚地看出。\Person{俄耳甫斯}与\Person{赫拉克勒斯}同时代；此外，据说归于他的所有作品都是由\Place{雅典}的\Person{奥诺玛克里托斯}所作，他生活在\Place{庇西特拉图}家族统治时期，大约第五十届奥林匹克运动会。\Person{穆赛俄斯}是\Person{俄耳甫斯}的弟子。\Person{安菲翁}比\Place{特洛伊}围城早两代，这使我们不必为那些想了解信息的人搜集关于他的更多细节。\Person{德摩多科斯}和\Person{斐弥俄斯}恰好生活在\Place{特洛伊}战争时期；前者与求婚者住在一起，后者与费阿刻斯人住在一起。\Person{塔米里斯}和\Person{菲拉蒙}比他们稍早一些。因此，关于他们各自在各种体裁中的作品、他们的时代以及记录，我们已经非常详尽且精确地写过了。但为了完成尚缺的部分，我将对那些被视为智者的人给出我的解释。\Person{米诺斯}，被认为在一切智慧、心智敏锐和立法能力方面超群出众，生活在\Person{林叩斯}时代，后者在\Person{达那俄斯}之后，是自\Person{伊那科斯}之后的第十一代。\Person{吕库古}，在\Place{特洛伊}陷落之后很久才出生，为\Place{拉刻代蒙}人立法。\Person{德拉古}被发现大约生活在第三十九届奥林匹克运动会时期，\Person{梭伦}大约在第四十六届，\Person{毕达哥拉斯}大约在第六十二届。我们已经证明奥林匹克运动会始于\Place{特洛伊}陷落后407年。这些事实得到证明后，我们将简要评论七贤的年龄。其中最年长的\Person{泰勒斯}大约生活在第五十届奥林匹克运动会时期；关于后来者我已经简要谈过。\par

% structure: p0092 source=h2 target=section reason=relative-heading-rank; base=h2

\section{第四十二章 关于作者的结束陈述}

这些事，啊，希腊人，我——\Person{塔提安}，蛮族哲学的门徒——为你们编撰的。我生于亚述人之地，先受教于你们的学说，后来才从事我现在所宣告的。从今以后，既然知道天主是谁，以及祂的作为是什么，我就把自己呈献给你们，准备好接受关于我学说的审查，同时坚定不移地持守那合乎天主的生活方式。\par

\SourceRecord{Translated by J.E. Ryland.；Ante-Nicene Fathers；Vol. 2.；Edited by Alexander Roberts, James Donaldson, and A. Cleveland Coxe.；Buffalo, NY: Christian Literature Publishing Co.,；1885.；<http://www.newadvent.org/fathers/0202.htm>.}
